% Source: https://artofproblemsolving.com/community/c3h2550626_seemingly_correct_solution_implies_easily_disprovable_corollary
% Stable identity: aops-topic-2550626-q01
\textbf{(a)} The argument starts from the statement to be proved and derives a true statement from it. That shows only ``claim $\Rightarrow$ true statement'', which proves nothing about the claim. Here the implication used, $n^5\equiv n\pmod{10}\Rightarrow n^5\equiv n\pmod5$, cannot be reversed: a congruence mod $5$ does not imply the congruence mod $10$. The same flawed pattern ``proves'' the false statement about $n^2$. (Also, the theorem meant is Fermat's \emph{little} theorem, not Fermat's Last Theorem.)

\textbf{(b)} We show $10\mid n^5-n$.
\begin{itemize}
\item By Fermat's little theorem, $n^5\equiv n\pmod5$ for every integer $n$, so $5\mid n^5-n$.
\item $n^5$ and $n$ have the same parity, so $2\mid n^5-n$.
\end{itemize}
Since $\gcd(2,5)=1$, both conditions together give $10\mid n^5-n$, i.e. $n^5$ and $n$ have the same last digit. $\blacksquare$

The two directions matter: $n^5\equiv n$ modulo both $2$ and $5$ \emph{implies} the congruence mod $10$.

\textbf{(c)} For $n=2$, $n^2=4$ ends in a different digit. (The congruence $n^2\equiv n$ holds mod $2$, but fails mod $5$, e.g. $2^2=4\not\equiv2\pmod 5$.)
